\documentclass[12pt]{article}
\usepackage[utf8]{inputenc}
\usepackage[T1]{fontenc}
\usepackage{hyperref}
\title{Urhobo Language Resources for K-12 Curriculum Development}
\author{Compiled by Claude Sonnet 5 (claude-sonnet-5) for Andr\'as Kornai}
\date{}
\begin{document}
\maketitle

\section{Introduction}
Urhobo (ISO 639-3 urh) is a Southwestern Edoid language of Delta State, Nigeria, with about 1.1 million speakers. It is taught as a subject in some schools. The search found bibliographic leads more than directly accessible materials. Several results were unreliable PDF mirrors, and I have left those out.

\section{Primary Educational Resources}
The search reported that the Urhobo Studies Association at Delta State University produced an approved 9-Year Basic Education Curriculum for Urhobo (Primary 1--3, Primary 4--6, JSS 1--3) and a scheme of work. I did not see a stable URL for these documents, so I do not list one. Omniglot notes that Urhobo is used in primary schools:
\begin{itemize}
\item Omniglot, Urhobo: \url{https://www.omniglot.com/writing/urhobo.htm}
\end{itemize}

\section{Linguistic and Reference Materials}
\begin{itemize}
\item Wikipedia, Urhobo language (archived copy). It lists dictionaries, including Ukere's \emph{Urhobo-English dictionary} (1986) and a dictionary from the Urhobo Language Institute: \url{https://webcf.waybackmachine.org/web/20220622235804/https://en.wikipedia.org/wiki/Urhobo_language}
\item Wikipedia, Southwestern Edoid languages (classification context): \url{https://en.wikipedia.org/wiki/Southwestern_Edoid_languages}
\end{itemize}

\section{Religious and Community Resources}
Omniglot mentions a Bible translation. No direct link to it was found in the search.

\section{Digital and Audio Resources}
None identified in the search.

\section{Curriculum Development Context}
Urhobo is stable and has a recognised basic-education curriculum. The next steps would be to obtain that curriculum and its scheme of work, and to contact the Urhobo Studies Programme at Delta State University and the Urhobo Language Institute. The search named Rose Aziza as director of that programme, but I give no contact details.
\end{document}